%% ============================================================
%% 分片 11：原文第 5 节「Topological Field Theory In Two Dimensions」（原文 glshort.tex 起于第 3264 行）
%% 由 tools/harvmac2tex.py 从 arXiv 源码（hep-th/0604151v3）自动生成。
%% 只译英文散文；公式、编号、\label、\cite 一律不改。
%% ============================================================

\section{二维拓扑场论}
\label{sec:tqfttwo}

\begin{keytakeaways}
  \begin{itemize}[leftmargin=1.4em,itemsep=2pt,topsep=2pt]
    \item \textbf{物理图像}：把第~\ref{sec:tqftfour}~节得到的四维扭化规范理论放到 $M=\Sigma\times C$ 上紧化（$C$ 是一张黎曼面），低能有效理论就是以 Hitchin 模空间 $\MH$ 为靶空间的二维拓扑 sigma 模型；$\MH$ 的超凯勒结构使它自动具有 $(4,4)$ 超对称。
    \item \textbf{数学工具}：四维扭化参数族 $\Bbb{CP}^1_g$ 与二维 $(2,2)$ 超对称所需的复结构对 $(J_+,J_-)$（即 $\Bbb{CP}^1_h\times\Bbb{CP}^1_h$ 中的一点）可以对齐，嵌入由 $w_+=-t$、$w_-=t^{-1}$ 给出，这里 $t$ 是四维扭化参数。
    \item \textbf{关键结论}：四维 $S$-对偶在二维表现为镜像对称——它交换 $A$-模型与 $B$-模型，从而把几何 Langlands 纲领表述为 Hitchin 模空间上膜范畴之间的等价。
    \item \textbf{边界与展望}：当 $J_-\ne\pm J_+$ 时所得模型既非 $A$-模型也非 $B$-模型，需要广义复几何才能把它约化为 $A$-模型；这正是本节后半与后续各节的出发点。
  \end{itemize}
\end{keytakeaways}

在第~\ref{sec:tqftfour}~节中，我们分析了 ${\cal N}=4$ 超对称杨-米尔斯理论在一个一般四维流形 $M$ 上的一个扭化版本。扭化保留了那些在子群 $Spin'(4)\subset Spin(4)\times Spin(6)$ 下不变的超对称。一个超对称生成元在 $Spin'(4)$ 下不变的条件是
% 式 (5.1)
\setcounter{equation}{0}
\begin{equation}\label{eq:tory}
\left(\Gamma_{\mu\nu}+\Gamma_{4+\mu,4+\nu}\right)\epsilon=0,
~~\mu,\nu=0,\dots,3.
\end{equation}
  此外，全部 16 个超对称生成元还满足一个十维手征条件，
% 式 (5.2)
\setcounter{equation}{1}
\begin{equation}\label{eq:nory}
\Gamma_0\Gamma_1\dots\Gamma_9\epsilon=\begin{cases}\epsilon&
~\text{in Lorentz signature}\\-i\epsilon&~\text{in Euclidean
signature.}\\\end{cases}
\end{equation}
 在扭化理论中，${\EUN}=4$ 超对称杨-米尔斯理论的六个标量场 $\phi_0,\dots,\phi_5$ 当中有四个被重新解释为一个 1-形式 $\phi=\sum_{\mu=0}^3\phi_\mu\,dx^\mu$。两个标量场 $\phi_4$ 与 $\phi_5$ 没有被扭化，它们在一个群 $SO(2)_{\cal R}$（我们在第~\ref{sec:twisting}~节中把它的生成元记作 ${\cal K}$）下转动。把费米子也一并计入，扭化理论的对称群是 $Spin(2)_{\cal R}$，即 $SO(2)_{\cal R}$ 的双重覆盖。

扭化理论是一种形式上的构造，因为扭化破坏酉性，并且只在欧几里得号差下才成立。\footnote{要在 Lorentz 号差下构造扭化理论，我们需要一个合适的同态 $Spin(1,3)\to Spin(6)$。由于 $Spin(6)$ 是紧的，这样的非平凡同态并不存在。若改把 $Spin(6)$ 换成 $Spin(1,5)$ 来绕开这一困难，则与费米子的耦合就不再是厄米的，能量也不再下有界。}然而，假定我们把时间方向分离出来，取 $M=\Bbb{R}\times W$，其中 $\Bbb{R}$ 参数化时间，$W$ 是一个三维流形，并且 $M$ 上的度量为乘积度量。那么时间不参与扭化，扭化理论在 Lorentz 号差或欧几里得号差下都说得通，并且是酉的、物理上合理的。此时 $M$ 的和乐群（holonomy group）约化为 $SO(3)$，未破缺的超对称只需在子群 $Spin'(3)\subset Spin'(4)$ 下不变；等价地说，我们只需对 $\mu,\nu=1,\dots,3$ 施加 \eqref{eq:tory}。共有四个未破缺的超对称生成元，它们对 $\mu,\nu=1,\dots 3$ 满足 \eqref{eq:tory}。有三个标量没有被扭化，即 $\phi_0$、$\phi_4$ 与 $\phi_5$，它们在一个 $Spin(3)_{\cal R}$ 对称下转动，该对称扩充了任何 $M$ 上都存在的 $Spin(2)_{\cal R}$ 对称。

让我们进一步特殊化到 $M=\Bbb{R}^{1,1}\times C$（或 $\Bbb{R}^2\times C$）的情形，其中 $C$ 是一张黎曼面。$M$ 的和乐群约化为 $SO(2)$，而 \eqref{eq:tory} 归结为单个条件
\begin{equation*}
\tag*{(5.3)}
\Gamma_{2367}\epsilon=\epsilon,\label{eq:ytytag}
\end{equation*}
留下一个八维的未破缺超对称空间。四个标量 $\phi_0,\phi_1,\phi_4$ 与 $\phi_5$ 没有被扭化，而转动未扭化标量的对称群此时扩充为 $Spin(4)_{\cal R}$。

事实上，二维中的结构恰好就是 $(4,4)$ 超对称的结构。若把 \eqref{eq:ytytag} 与 \eqref{eq:nory} 结合起来，我们得到（在 Lorentz 号差下）% 式 (5.4)
\setcounter{equation}{3}
\begin{equation}\label{eq:byty}
\Gamma_{01}\epsilon
=\Gamma_{4589}\epsilon.
\end{equation}
 这里 $\Gamma_{01}$ 是区分 $\Bbb{R}^{1,1}$ 的 Lorentz 群 $SO(1,1)$ 的两个旋量表示的算子。类似地，$\Gamma_{4589}$ 区分 $Spin(4)_{\cal R}=SU(2)'\times SU(2)''$ 的两个旋量表示。这两个旋量表示在 $SU(2)'\times SU(2)''$ 下分别按 $({\bf 2},{\bf 1})$ 与 $({\bf 1},{\bf 2})$ 变换。式 \eqref{eq:byty} 说明：对所有的二维超对称，$\Gamma_{01}$ 与 $\Gamma_{4589}$ 的本征值都相等。

于是，在适当选取 $\Bbb{R}^{1,1}$ 的定向之后，超对称按 $({\bf 2},{\bf 1})^-\oplus({\bf 1},{\bf 2})^+$ 变换，其中上标标记 $\Gamma_{01}$ 的本征值。但这正是二维中 $(4,4)$ 超对称的通常结构：一个 $SU(2)$ 的 $R$-对称群——即 $SU(2)'$——作用在负 $Spin(1,1)$ 手征的超对称上，而第二个这样的群——$SU(2)''$——作用在正手征的超对称上。事实上，在第~\ref{sec:compacthitchin}~节中，我们已经把低能有效理论确定为靶空间是超凯勒流形 $\MH$ 的 sigma 模型。一般来说，靶空间为超凯勒流形的 sigma 模型具有 $(4,4)$ 超对称 \cite{ref65}。

除了 $Spin(4)_{\cal R}$ 整体对称之外，$\Bbb{R}^{1,1}\times C$ 上的扭化规范理论还有一个额外对称，它之所以出现，是因为我们据以扭化的 $C$ 的和乐群是阿贝尔的。即使在扭化 $\phi_2$ 与 $\phi_3$ 之后，仍然可以对 $\phi_2-\phi_3$ 平面作转动。这就是第~\ref{sec:mhhyper}~节中称为 ${\cal U}_1$ 的对称群，它作用在 Hitchin 对 $(E,\varphi)$ 上，把 $\varphi$ 的相位转动。

\begin{yizhu}[原文疑误]
此处原文说的是把式 (5.3) 与式 (5.2) 结合起来。harvmac 源码里 $\Gamma_{2367}\epsilon=\epsilon$ 这个词条被写成了无编号的 \texttt{tag*} 公式（源文件用同名宏 \verb|\yty| 定义了两次，第二次即 (5.3)），译本已为它补上标签 \texttt{eq:ytytag}，故 \eqref{eq:ytytag} 正确指向 (5.3)。按：这是排版脚手架问题，不是原文笔误。
\end{yizhu}

\begin{yizhu}[易错点]
扭化只在欧几里得号差下成立：它破坏酉性（脚注给出了理由——紧群 $Spin(6)$ 与非紧的 $Spin(1,3)$ 之间不存在非平凡同态）。但一旦把时间方向分离出来写成 $M=\Bbb{R}\times W$，扭化就与时间无关，理论在 Lorentz 号差下也可以既酉又物理。读本文时不要因为第一句话就把扭化理论当作纯形式的玩具。
\end{yizhu}

\begin{yizhu}[含糊点拨]
注意未被扭化的标量数目随流形而变：一般四维流形 $M$ 上为 2 个（$\phi_4,\phi_5$），$\Bbb{R}\times W$ 上为 3 个（$\phi_0,\phi_4,\phi_5$），$\Bbb{R}^{1,1}\times C$ 上为 4 个（$\phi_0,\phi_1,\phi_4,\phi_5$）；相应地 $R$-对称从 $Spin(2)_{\cal R}$ 逐步扩充为 $Spin(3)_{\cal R}$、$Spin(4)_{\cal R}$。按：这些数目正是二维 $(4,4)$ 超对称中每个「4」的来源。
\end{yizhu}

\begin{yizhu}[物理图像]
$\Gamma_{01}\epsilon=\Gamma_{4589}\epsilon$ 把两件不同的事绑在一起：$\Gamma_{01}$ 标记 $\Bbb{R}^{1,1}$ 上 $SO(1,1)$ 的两种手征，$\Gamma_{4589}$ 标记 $R$-对称群 $SU(2)'\times SU(2)''$ 的两个二重态。二者本征值相等，意味着二维的手征与 $R$-对称的 $SU(2)$ 标记同步，这正是二维 $(4,4)$ 超对称的标准结构；也正因为如此，靶空间为超凯勒流形的 sigma 模型才「自动」是 $(4,4)$ 的。
\end{yizhu}

\begin{yizhu}[含糊点拨]
本节最后一段的 ${\cal U}_1$ 对称值得留意：它之所以存在，是因为用来扭化的 $C$ 的和乐群 $SO(2)$ 是阿贝尔的——扭化用掉了 $\phi_2,\phi_3$ 所在的二维平面，但该平面内仍可整体转动。这个 ${\cal U}_1$ 与复结构 $I_w$ 族上的 $\Bbb{C}^*$ 作用（第~\ref{sec:complexstr}~节末）是同一件事在 $|\lambda|=1$ 时的限制，后面讨论 T-对偶与镜像对称时还会用到。
\end{yizhu}
